\section{Related Work}

\subsection{Remote Physiological Sensing via Webcam}
The past decade has witnessed a revolution in non-contact physiological sensing, driven by the ubiquity of high-resolution webcams and advances in computer vision. The most prominent technique, remote photoplethysmography (rPPG), extracts the blood volume pulse (BVP) signal by analyzing subtle chromatic variations in the skin caused by cardiovascular activity. Since the seminal work of Verkruysse et al., who laid the optical foundations for rPPG using ambient light, the field has rapidly evolved from signal processing approaches (e.g., ICA, POS \cite{wang_pos}) to deep learning-based methods (e.g., DeepPhys, PhysNet).

McDuff et al. \cite{mcduff_rppg} provide a comprehensive review of these advances, highlighting that modern rPPG algorithms can achieve heart rate estimation accuracy comparable to contact-based pulse oximeters ($MAE < 3$ bpm) under stationary conditions. Similarly, webcam-based eye tracking has matured significantly. WebGazer.js \cite{webgazer} demonstrated that a self-calibrating regression model running entirely in the browser could achieve gaze estimation with an accuracy of approx. $4^{\circ}$, sufficient for attention tracking on standard screens.

However, a critical gap remains between \textit{signal extraction} and \textit{affective interpretation}. The vast majority of rPPG studies validate their methods on labeled datasets (e.g., UBFC-rPPG, MMSE-HR) where the ground truth is the physiological signal itself (HR, HRV). Few studies venture into the downstream task of interpreting these signals in complex, unconstrained environments. When they do, they often rely on simplistic arousal models (e.g., assuming heart rate variability correlates linearly with cognitive load), neglecting the complex interplay between physiology and context.

\subsection{Context-Aware Affective Computing}
The necessity of context in emotion recognition is well-acknowledged in theoretical frameworks but remains difficult to implement in practice. Ma and Hovy \cite{context_aware_ac} identify "multimodal context"—including location, social setting, and temporal dynamics—as essential for disambiguating affective signals. Without context, a smile could indicate joy (Duchenne smile) or social politeness; a furrowed brow could indicate anger or concentrated attention.

Traditional approaches to context-awareness involve adding sensor modalities (e.g., GPS, accelerometer) to classify the user's activity (e.g., "Walking," "Sitting"). However, this physical context does not fully capture the \textit{psychological context}. The Biopsychosocial Model \cite{biopsychosocial_model} and Polyvagal Theory \cite{polyvagal} argue that the key contextual variable is "Safety." Porges \cite{porges_safe_haven} introduced the concept of "neuroception"—a subconscious neural process that scans the environment for cues of safety, danger, or life-threat.

When an environment is neurocepted as "Safe" (e.g., a quiet home), the ventral vagal complex is active, promoting social engagement and dampening sympathetic reactivity. When it is neurocepted as "Danger" (e.g., a formal lab with evaluative observers), the sympathetic nervous system is disinhibited (the "Safe Haven" effect is removed), leading to hyper-reactivity. Our work specifically investigates this "Safety Context," which is often invisible to GPS or activity classifiers but fundamental to affective processing.

\subsection{Single-Case Experimental Designs (SCED)}
While large-$N$ randomized controlled trials (RCTs) are the gold standard for population-level generalization (e.g., testing a new drug), they are often ill-suited for studying personalized, dynamic phenomena where inter-individual variability is high. In Affective Computing, basal physiological levels (e.g., resting Heart Rate) vary wildly between individuals, creating noise that masks subtle experimental effects in group designs.

Single-Case Experimental Design (SCED), also known as N-of-1 trials, offers a rigorous alternative. Originating in behavioral psychology (Skinner) and clinical medicine \cite{barlow_sced}, SCED involves the intensive, prospective measurement of a single subject over time, using randomized phases (e.g., A-B-A-B designs) to establish causal relationships within that individual.

Kazemitabar et al. \cite{kazemitabar_sced_hci} argue for the utility of SCED in Human-Computer Interaction (HCI) to understand individual variability in technology usage. By treating the individual as their own control, SCED achieves high internal validity and can detect small effect sizes that would be washed out in group averages. This makes it an ideal framework for developing "Personalized Affective Computing" models, which must learn the specific physiological signature of a single user before generalizing to others.
